% !TeX program = lualatex
% !TeX encoding = utf8
% !TeX spellcheck = uk_UA
% !TeX root = ../ElMagConspect.tex

%=========================================================
\chapter{Система рівнянь Максвелла}\label{\currfilebase}
\graphicspath{{\currfilebase/Pictures}}
%=========================================================
\epigraph{Кількісне, математичне формулювання законів поля дано в так званих рівняннях Максвелла. Експериментальні факти призвели до
	формулювання цих рівнянь, але зміст їх значно багатший [\ldots]. Їхня проста форма приховує глибину, що виявляється тільки при ретельному
	вивченні.
}{А. Ейнштейн, Л. Інфельд, \enquote{Еволюція фізики}}

\section*{Основоположники теорії електромагнітного поля}

Теорію електромагнітного поля, початки якої заклав \emphx{Фарадей}, математично завершив \emphx{Максвелл}. Однією з найважливіших нових ідей,
висунутих Максвеллом, була думка про \emphx{симетрію} у взаємозалежності електричного й магнітного полів: якщо змінне магнітне поле породжує
електричне (закон Фарадея, гл.~\ref{ElMagInduction}), то, за красивою симетрією, змінне електричне поле мусить породжувати магнітне. Саме ця ідея,
як ми зараз побачимо, і привела Максвелла до відкриття \emphx{струму зміщення} --- останньої відсутньої ланки повної системи рівнянь
електродинаміки.

\begin{figure}[h!]
	\centering
	\begin{minipage}[t]{0.45\linewidth}\centering
		\includegraphics[height=4cm]{Faraday}\\
		\footnotesize Майкл Фарадей (1791--1867)
	\end{minipage}
    \hfill
	\begin{minipage}[t]{0.45\linewidth}\centering
		\includegraphics[height=4cm]{james-clerk-maxwell}\\
		\footnotesize Джеймс Клерк Максвелл (1831--1879)
	\end{minipage}
	\caption{Основоположники класичної теорії електромагнітного поля\label{pic:FaradayMaxwell}}
\end{figure}

Формулювання рівнянь Максвелла А.~Ейнштейн і Л.~Інфельд у книзі \enquote{Еволюція фізики} назвали \enquote{найважливішою подією з часу Ньютона}
--- не лише через цінність їхнього змісту, а й тому, що вони дають зразок \emphx{нового типу законів}: \emphx{рівняння Максвелла суть закони, що
	виражають структуру поля}, а не закони руху окремих частинок.

%% --------------------------------------------------------
\section{Струм зміщення}
%% --------------------------------------------------------

\subsection*{Протиріччя в законах магнетизму}

Теорема про циркуляцію для постійного магнітного поля, яку ми встановили в гл.~\ref{MagFieldMedia},
\begin{equation}\label{eq:CirculationHVacuumOld}
	\oint\limits_L \Hfield\cdot d\vect r = \frac{4\pi}{c}\iint\limits_S \vect j\cdot d\vect S, \qquad \rot\Hfield = \frac{4\pi}{c}\vect j,
\end{equation}
виявляється \emphx{неправильною} у випадку \emphx{змінного} електричного поля (рис.~\ref{pic:CircThm}). Застосуємо операцію $\divg$ до
диференціальної форми цього рівняння: оскільки дивергенція будь-якого ротора тотожно дорівнює нулю, $\divg\rot\Hfield \equiv 0$, отримуємо
\begin{equation}\label{eq:DivJZero}
	\divg\vect j = 0.
\end{equation}

\begin{SCfigure}[1][h!]\centering
	\localinput{circ_thm.tikz}
	\caption{Контур $L$, натягнута на нього поверхня $S$ та провідник зі струмом, що її
		пронизує\label{pic:CircThm}}
\end{SCfigure}

Але, з іншого боку, за законом збереження заряду (рівняння неперервності, гл.~\ref{SteadyCurrent}), якщо густина заряду змінюється з часом,
$\partial\rho/\partial t \neq 0$, то
\begin{equation}\label{eq:ContinuityViolation}
	\divg\vect j = -\frac{\partial\rho}{\partial t} \neq 0.
\end{equation}
Рівняння~\eqref{eq:DivJZero} і~\eqref{eq:ContinuityViolation} відверто \emphx{суперечать} одне одному щоразу, коли $\partial\rho/\partial t\neq0$
--- а отже, теорему про циркуляцію~\eqref{eq:CirculationHVacuumOld} необхідно \emphx{видозмінити}.

\subsection*{Гіпотеза Максвелла}

Щоб усунути цю суперечність, Дж.~Максвелл увів поняття \emphx{струму зміщення} $\vect j_\text{зм}$, доповнивши теорему про циркуляцію так, щоб
закон збереження заряду виконувався за будь-яких обставин:
\begin{equation}\label{eq:DisplacementCurrentDef}
	\rot\Hfield = \frac{4\pi}{c}\left(\vect j+\vect j_\text{зм}\right).
\end{equation}
Застосовуючи операцію $\divg$ до цього рівняння:
\begin{equation*}
	\divg\left(\vect j+\vect j_\text{зм}\right) = 0 \quad\Rightarrow\quad \divg\vect j_\text{зм} = \frac{\partial\rho}{\partial t}
\end{equation*}
(тут ми скористались рівнянням~\eqref{eq:ContinuityViolation}). За теоремою Гаусса для електричного поля, $\rho = \dfrac1{4\pi}\divg\Dfield$
(гл.~\ref{FileldInDielectrics}), тому
\begin{equation*}
	\divg\vect j_\text{зм} = \frac{\partial}{\partial t}\left(\frac1{4\pi}\divg\Dfield\right) = \divg\left(\frac1{4\pi}\frac{\partial\Dfield}{\partial
	t}\right).
\end{equation*}
Найпростіше (і, як показує подальша перевірка експериментом, єдино правильне) розв'язання цього рівняння --- покласти
\begin{equation}\label{eq:DisplacementCurrentDeriv}
	\tcbhighmath{
		\vect j_\text{зм} = \frac1{4\pi}\frac{\partial\Dfield}{\partial t}.
	}
\end{equation}

\begin{Attention}
	Рівняння~$\divg\vect j_\text{зм} = \divg\left(\tfrac1{4\pi}\partial\Dfield/\partial t\right)$ саме собою \emphx{не доводить}, що
	$\vect j_\text{зм}$ дорівнює точно $\tfrac1{4\pi}\partial\Dfield/\partial t$: до цього виразу можна було б, узагалі кажучи, додати довільне
	бездивергентне поле ($\divg\vect f=0$). Максвелл обрав \emphx{найпростіший} можливий варіант --- без жодного додаткового доданка; те, що саме
	цей вибір відповідає дійсності (а не просто математичній зручності), --- це вже суто \emphx{експериментальний} факт, підтверджений, зокрема,
	дослідами Г.~Герца з випромінювання електромагнітних хвиль (див. закінчення цього розділу).
\end{Attention}

\subsection*{Теорема про циркуляцію магнітного поля з урахуванням струму зміщення}

Отже, теорема про циркуляцію для магнітного поля, узгоджена із законом збереження заряду, набуває остаточного вигляду:
\begin{equation}\label{eq:CirculationHFull}
	\tcbhighmath{
		\rot\Hfield = \frac{4\pi}{c}\vect j + \frac1c\frac{\partial\Dfield}{\partial t}.
	}
\end{equation}
В інтегральній формі:
\begin{equation}\label{eq:CirculationHFullInt}
	\tcbhighmath{
		\oint\limits_L \Hfield\cdot d\vect r = \frac{4\pi}{c}\iint\limits_S \vect j\cdot d\vect S + \frac1c\iint\limits_S
		\frac{\partial\Dfield}{\partial t}\cdot d\vect S,
	}
\end{equation}
де
\begin{equation}\label{eq:DisplacementCurrentIntegral}
	I_\text{зм} = \frac1{4\pi}\iint\limits_S \frac{\partial\Dfield}{\partial t}\cdot d\vect S
\end{equation}
--- струм зміщення, що пронизує поверхню $S$, натягнуту на контур $L$. Отже, згідно з гіпотезою Максвелла, \emphx{змінне електричне поле поряд зі
	звичайними струмами провідності також створює магнітне поле}.

\subsection*{Порівняння закону Фарадея та гіпотези Максвелла}

Порівняємо закон електромагнітної індукції Фарадея (гл.~\ref{ElMagInduction}) і \enquote{оновлену} теорему про циркуляцію за відсутності струмів
провідності ($\vect j=0$) у вакуумі ($\epsilon=\mu=1$, тобто $\Efield=\Dfield$, $\Hfield=\Bfield$):
\begin{equation}\label{eq:FaradayVsMagnetoelectric}
	\underbrace{\rot\Efield = -\frac1c\frac{\partial\Bfield}{\partial t}}_\text{закон електромагнітної індукції} \qquad\qquad
	\underbrace{\rot\Bfield = +\frac1c\frac{\partial\Efield}{\partial t}}_\text{\enquote{закон магнітоелектричної індукції}}.
\end{equation}

\begin{Attention}
	Ця пара рівнянь --- вершина ідеї Максвелла про \emphx{симетрію} електрики й магнетизму: \emphx{змінне в часі магнітне поле породжує вихрове
		електричне поле, а змінне в часі електричне поле породжує магнітне поле}. Єдина відмінність --- знак: у законі Фарадея \enquote{мінус}
	(правило Ленца), у гіпотезі Максвелла --- \enquote{плюс}. Ця відносно невелика асиметрія знаків, утім, має глибокі наслідки: саме вона (разом з
	рівняннями $\divg\Efield=0$, $\divg\Bfield=0$ у вакуумі) приводить, як буде показано в наступних розділах курсу, до існування
	\emphx{електромагнітних хвиль}, що біжать зі швидкістю світла $c$, а не до експоненційно загасаючих чи зростаючих розв'язків.
\end{Attention}

\subsection*{Приклад 1: радіальне стікання заряду з кулі}

Нехай куля несе заряд $Q(t)$, що стікає в зовнішнє (слабко провідне) середовище рівномірно в усіх напрямках (рис.~\ref{pic:RadialLeak}). Чи
виникає при цьому магнітне поле?

\begin{SCfigure}[2][h!]\centering
	\localinput{radial-leak.tikz}
	\caption{Радіальне стікання заряду з кулі: струм провідності $I$ і струм зміщення
		$I_\text{зм}$\label{pic:RadialLeak}}
\end{SCfigure}

Стікання заряду створює радіальний струм провідності $I = -\dfrac{\partial Q}{\partial t}$. Поле кулі на відстані $r$, за теоремою Гаусса,
$\Efield(r) = \dfrac{Q(t)}{r^3}\vect r$ (і теж спадає з часом), тому крізь будь-яку сферу радіуса $r>R$ додатково тече струм зміщення:
\begin{equation}\label{eq:RadialLeakCurrent}
	I_\text{зм} = \frac1{4\pi}\oiint\limits_S \frac{\partial\Efield}{\partial t}\cdot d\vect S = \frac{\partial Q}{\partial t} = -I.
\end{equation}
За теоремою про циркуляцію, застосованою до довільного кола $L$, що оточує напрям витікання струму (з відповідною поверхнею $S$, що її
перетинають і провідність, і струм зміщення):
\begin{equation}\label{eq:RadialLeakZeroB}
	\oint\limits_L \Hfield\cdot d\vect r = \frac{4\pi}c\left(I+I_\text{зм}\right) = 0 \quad\Rightarrow\quad \Hfield=\Bfield=0.
\end{equation}

\begin{Attention}
	У цій задачі, попри те, що заряд явно \emphx{рухається} (тече струм!), магнітне поле \emphx{не виникає взагалі}: струм зміщення в кожній
	точці простору точно компенсує струм провідності. Це наочно ілюструє, що струм зміщення --- не якась математична поправка \enquote{про око},
	а реальний фізичний внесок у джерела магнітного поля, необхідний для узгодженості всієї теорії.
\end{Attention}

\subsection*{Приклад 2: струм зміщення в конденсаторі}

Розглянемо коло з плоским конденсатором, яке заряджається від джерела струмом $I$ (рис.~\ref{pic:TwoSurfaces}). Візьмемо контур $L$, що охоплює
провід поблизу конденсатора, і натягнемо на нього \emphx{дві різні} поверхні: одну (ліворуч) --- що перетинає сам провід, і другу (праворуч) ---
що проходить між обкладками конденсатора, не перетинаючи жодного проводу.

\begin{figure}[h!]\centering
	\localinput{two-surfaces-loop.tikz}
	\caption{Дві різні поверхні, натягнуті на один і той самий контур $L$\label{pic:TwoSurfaces}}
\end{figure}

Для лівої поверхні циркуляція визначається звичайним струмом провідності: $\oint_L\Hfield\cdot d\vect r = \tfrac{4\pi}{c}I$. Для правої
поверхні провідного струму немає взагалі (вона проходить у проміжку між обкладками, де є лише електричне поле, що змінюється в часі), тому
циркуляція визначається \emphx{лише} струмом зміщення: $\oint_L\Hfield\cdot d\vect r = \tfrac{4\pi}c I_\text{зм}$.

\begin{equation}\label{eq:CapacitorTwoSurfaces}
	\tcbhighmath{
		I = I_\text{зм}
	}
\end{equation}

Оскільки циркуляція вектора $\Hfield$ по тому самому контуру $L$ не може залежати від вибору поверхні, натягнутої на цей контур (це --- пряма
математична властивість теореми Стокса), звідси випливає $I=I_\text{зм}$: \emphx{всередині конденсатора \enquote{протікає} струм зміщення, який
	замикає коло}. Розглянемо докладніше, яке магнітне поле він створює (рис.~\ref{pic:CapacitorChargingField}).

\begin{figure}[h!]\centering
	\localinput{capacitor-charging-field.tikz}
	\caption{Магнітне поле, створене струмом зміщення всередині конденсатора під час його
		зарядки\label{pic:CapacitorChargingField}}
\end{figure}

\begin{Attention}
	\emphx{Струми зміщення замикають струми провідності і створюють магнітне поле точнісінько так само, як і звичайні струми провідності}. Проте
	вони мають і принципову відмінність: струм зміщення \emphx{не} пов'язаний з тепловиділенням --- до нього незастосовні ні закон Ома, ні закон
	Джоуля--Ленца (гл.~\ref{SteadyCurrent}), адже він не пов'язаний з переміщенням реальних носіїв заряду проти сил опору, а лише зі зміною
	електричного поля в просторі.
\end{Attention}

%% --------------------------------------------------------
\section{Система рівнянь Максвелла}
%% --------------------------------------------------------

Доповнивши основні закони електромагнетизму, встановлені в попередніх розділах, гіпотезою про струм зміщення, Максвелл отримав повну систему
фундаментальних рівнянь електродинаміки. Таких рівнянь --- чотири (табл.~\ref{tab:MaxwellSystem}).

\begin{table}[h!]\small\centering
	\caption{Система рівнянь Максвелла\label{tab:MaxwellSystem}}
		\begin{tblr}%
			{
			colspec={|Q[l,m, 4cm, bg=yellow!10, font=\small]|[1pt, white]X[l,m, bg=red!10]|[1pt, white]Q[l,m, 4cm, bg=blue!10]|},
			row{1}={c, fg=white, bg=themecolorlight, font=\bfseries\small},
			cell{2-5}{2-Z}={l, m, mode=dmath, font=\small},
			%hlines,
			}
			\hline
			Рівняння                                                                 & Інтегральна форма                                                                       & Диференціальна форма \\
			\hline\hline
			Теорема Гаусса для магнітного поля                                       & \oiint\limits_S \Bfield\cdot d\vect{S} =  0                                             & \Div\Bfield = 0      \\
			\hline
			Теорема про циркуляцію для електричного поля                             & \oint\limits_L \Efield\cdot d\vect{r} = -
			\frac1c\iint\limits_S \frac{\partial\Bfield}{\partial t} \cdot d\vect{S} & \Rot\Efield = - \frac1c \frac{\partial\Bfield}{\partial t}                                                     \\
			\hline\hline
			Теорема Гаусса для електричного поля                                     & \oiint\limits_S \Dfield\cdot d\vect{S} = 4\pi \iiint\limits_{V} \rho dV                 & \Div\Dfield =
			4\pi\rho                                                                                                                                                                                  \\
			\hline
			Теорема про циркуляцію для магнітного поля                               & \oint\limits_L \Hfield\cdot d\vect{r} = \frac{4\pi}{c} \iint\limits_S \left( \vect{j} +
			\frac1{4\pi}\frac{\partial\Dfield}{\partial t} \right) \cdot d\vect{S}   & \Rot\Hfield = \frac{4\pi}{c}\vect{j} + \frac1c
			\frac{\partial\Dfield}{\partial
				t}
			\\
			\hline
		\end{tblr}
\end{table}

\subsection*{Вираз електричного поля через потенціали}

Рівняння Максвелла природно групуються парами: \emphx{перша} пара (теорема Гаусса для $\Bfield$ і циркуляція $\Efield$) --- рівняння \emphx{без}
явних джерел (зарядів і струмів), \emphx{друга} пара --- рівняння \emphx{з} джерелами.

Теорема Гаусса для $\Bfield$ дає змогу, як і в магнітостатиці (гл.~\ref{MagFieldVacuum}), ввести вектор-потенціал:
\begin{equation*}
	\Bfield = \rot\vect A.
\end{equation*}
Підставляючи це в закон індукції, $\rot\Efield = -\dfrac1c\parttime{\Bfield} = -\dfrac1c\parttime{}(\rot\vect A) = -\rot\left(\dfrac1c
\parttime{\vect A}\right)$, отримуємо
\begin{equation*}
	\rot\left(\Efield+\frac1c\parttime{\vect A}\right) = 0.
\end{equation*}
Оскільки ротор цього виразу тотожно дорівнює нулю, сам вираз можна подати як градієнт деякої скалярної функції $-\phi$:
\begin{equation}\label{eq:EFromPotentials}
	\tcbhighmath{
		\Efield = -\grad\phi - \frac1c\frac{\partial\vect A}{\partial t}.
	}
\end{equation}
У випадку сталих у часі полів $\parttime{\vect A}=0$, і $\Efield=-\grad\phi$ --- введена тут функція $\phi$ природно збігається зі звичним
скалярним потенціалом електростатики.

\begin{Attention}
	Формула~\eqref{eq:EFromPotentials} показує, що в загальному, нестаціонарному випадку електричне поле визначається \emphx{не тільки}
	скалярним потенціалом $\phi$ (як в електростатиці), а й швидкістю зміни вектор-потенціала $\vect A$: частина $\Efield$ породжується не
	зарядами напряму, а через посередництво змінного магнітного поля --- це і є, по суті, вихрова, \enquote{індукційна} частина електричного
	поля, з якою ми вже зустрічались у гл.~\ref{ElMagInduction}.
\end{Attention}

\subsection*{Стаціонарні поля}

У стаціонарному випадку всі часткові похідні за часом дорівнюють нулю, $\parttime{}=0$, і повна система рівнянь Максвелла розпадається на дві
\emphx{незалежні} системи --- для електростатики й магнітостатики (рис.~\ref{pic:StaticsComparison}):
\begin{equation}\label{eq:StaticDecoupling}
	\underbrace{\begin{cases}
			\divg\Dfield = 4\pi\rho, \\
			\rot\Efield = 0
		\end{cases}}_\text{електростатика} \qquad\qquad
	\underbrace{\begin{cases}
			\divg\Bfield = 0, \\
			\rot\Hfield = \dfrac{4\pi}c\vect j
		\end{cases}}_\text{магнітостатика}.
\end{equation}

\begin{figure}[h!]\centering
	\localinput{statics-comparison.tikz}
	\caption{У стаціонарному випадку електричне й магнітне поля \enquote{не знають} одне про
		одного\label{pic:StaticsComparison}}
\end{figure}

\begin{Attention}
	Це --- глибокий концептуальний висновок: усі розділи курсу, які передували цьому (електростатика, магнітостатика), --- не окремі теорії, а
	лише \emphx{частинні випадки} єдиної теорії електромагнітного поля Максвелла, що виникають у граничному випадку, коли поля не залежать від
	часу. \emphx{Зв'язок} електрики й магнетизму (взаємне породження одне одного) --- принципово \emphx{нестаціонарне} явище, яке зникає, щойно
	всі процеси застигають у часі.
\end{Attention}

\subsection*{Матеріальні рівняння}

Рівняння Максвелла --- це рівняння \emphx{загальної} структури поля; щоб розв'язати конкретну задачу, їх необхідно доповнити
\emphx{матеріальними рівняннями}, що пов'язують поля $\Dfield$ та $\Efield$, з одного боку, і $\Hfield$ та $\Bfield$ --- з іншого, і залежать від
властивостей конкретного середовища (табл.~\ref{tab:ConstitutiveRelations}).

\begin{table}[h!]\small\centering
	\caption{Матеріальні рівняння\label{tab:ConstitutiveRelations}}
		\begin{tblr}%
			{
			colspec={|X[l,m, bg=yellow!10, font=\small]|[1pt, white]Q[l,m, 0.2\linewidth,bg=red!10]|},
			row{1}={c, fg=white, bg=themecolorlight, font=\bfseries},
			cell{2-Z}{2-Z}={l, m, mode=dmath, font=\small},
			hlines,
			}
			                                                                                           & Вираз                            \\
			Вектор індукції електричного поля                                                          & \Dfield = \Efield + 4\pi\vect{P} \\
			Вектор поляризації (для лінійних ізотропних речовин)                                       & \vect{P} = \chi_e\Efield         \\
			Діелектрична проникніть  (для лінійних ізотропних речовин)                                 & \epsilon = 1 + 4\pi \chi_e       \\
			Зв'язок $\Dfield$ та $\Efield$ (для лінійних ізотропних речовин)                           & \Dfield = \epsilon\Efield        \\
			\hline
			Вектор напруженості магнітного поля                                                        & \Hfield = \Bfield - 4\pi\vect{J} \\
			Вектор намагнічування (для лінійних ізотропних речовин)                                    & \vect{J} = \chi_e\Hfield         \\
			Магнітна проникніть  (для лінійних ізотропних речовин)                                     & \mu = 1 + 4\pi \chi_m            \\
			Зв'язок $\Hfield$ та $\Bfield$ (для лінійних ізотропних речовин)                           & \Bfield = \mu\Hfield             \\
			\hline
			У разі струму, спричиненого електричним полем у провідному середовищі, має місце закон Ома & \vect{j} = \lambda\Efield
		\end{tblr}
\end{table}

\subsection*{Граничні умови}

Диференціальні рівняння Максвелла необхідно доповнити граничними умовами на межі розділу двох середовищ; ці умови неявно містяться в
\emphx{інтегральній} формі рівнянь Максвелла і виводяться з них тим самим способом (циліндр Гаусса, контур циркуляції), що й у попередніх
розділах курсу. Граничні умови для змінних полів мають \emphx{той самий вигляд}, що й для стаціонарних (табл.~\ref{tab:BoundaryConditionsFull}).

\begin{table}[h!]\small\centering
	\caption{Граничні умови для електромагнітного поля\label{tab:BoundaryConditionsFull}}
		\begin{tblr}%
			{
			colspec={|X[l,m, bg=yellow!10, font=\small]|[1pt, white]X[c,m, bg=red!10]|},
			row{1}={c, fg=white, bg=themecolorlight, font=\bfseries},
			cell{2-Z}{1-Z}={c, m, mode=dmath, font=},
			%hlines,
			}
			\hline
			Умови для електричних векторів & Умови для магнітних векторів \\
			\hline
			D_{2n} - D_{1n} = 4\pi\sigma
			                               &
			B_{1n} = B_{2n}
			\\
			\hline
			\hline
			E_{1\tau} = E_{2\tau}
			                               &
			\left[\vect{n}\times\Hfield_2\right] - \left[\vect{n}\times\Hfield_1\right] = \frac{4\pi}{c}\vect{i}
			\\
			\hline
		\end{tblr}
\end{table}

Тут $\sigma$ --- поверхнева густина \emphx{вільних} електричних зарядів, а $\vect i$ --- поверхнева густина струму провідності на межі розділу.
Якщо поверхневих струмів немає, гранична умова для тангенціальної компоненти $\Hfield$ спрощується:
\begin{equation*}
	H_{1\tau} = H_{2\tau}.
\end{equation*}

%% --------------------------------------------------------
\section{Енергія електромагнітного поля}
%% --------------------------------------------------------

\subsection*{Закон збереження енергії електромагнітного поля}

Нехай у деякій області простору присутнє електромагнітне поле й заряджені частинки, що з ним взаємодіють (рис.~\ref{pic:EnergyFlowRegion}). Якщо
ця система \emphx{замкнена}, її енергія зберігається; якщо ж \emphx{незамкнена}, енергія може як \emphx{втікати} в цю область, так і
\emphx{витікати} з неї. Якісно закон збереження енергії для такої області можна записати так:
\begin{equation*}
	(\text{зміна енергії в об'ємі}) = (\text{рух частинок}) + (\text{потік енергії назовні}),
\end{equation*}
або, виражаючи цікаву для нас величину --- енергію, витрачену на рух заряджених частинок:
\begin{equation*}
	(\text{рух частинок}) = -(\text{зміна енергії в об'ємі}) - (\text{потік енергії назовні}).
\end{equation*}

\begin{SCfigure}[2][h!]\centering
	\localinput{energy-flow-region.tikz}
	\caption{Заряджені частинки в області, обмеженій замкненою поверхнею: енергія
		електромагнітного поля може втікати й витікати крізь цю поверхню\label{pic:EnergyFlowRegion}}
\end{SCfigure}

Рухати заряджені частинки, як ми твердо встановили ще в гл.~\ref{MagFieldVacuum}, може \emphx{лише електричне} поле (магнітне поле роботи не
виконує). Енергія, яку електричне поле витрачає на розгін частинок в одиниці об'єму за одиницю часу, --- це якраз джоулева потужність
$\vect j\cdot\Efield$ (гл.~\ref{SteadyCurrent}).

\subsection*{Виведення закону збереження}

Запишемо закон Максвелла для $\rot\Hfield$~\eqref{eq:CirculationHFull}, розв'язавши його відносно $\vect j$:
\begin{equation*}
	\vect j = \frac{c}{4\pi}\rot\Hfield - \frac1c\parttime{\Dfield}.
\end{equation*}
Помножимо скалярно обидві частини на $\Efield$:
\begin{equation*}
	\vect j\cdot\Efield = \frac{c}{4\pi}\Efield\cdot\rot\Hfield - \frac1{4\pi}\Efield\cdot\parttime{\Dfield}.
\end{equation*}
Скористаємось тотожністю векторного аналізу $\Efield\cdot\rot\Hfield = -\divg\left[\Efield\times\Hfield\right] + \Hfield\cdot\rot\Efield$ та
матеріальним рівнянням $\Dfield=\epsilon\Efield$ (так що $\Efield\cdot\parttime{\Dfield} = \epsilon\Efield\cdot\parttime{\Efield} =
\parttime{}\left(\tfrac12\epsilon E^2\right)$):
\begin{equation*}
	\vect j\cdot\Efield = -\divg\left(\frac{c}{4\pi}\left[\Efield\times\Hfield\right]\right) + \frac{c}{4\pi}\Hfield\cdot\rot\Efield -
	\parttime{}\left(\frac1{8\pi}\epsilon E^2\right).
\end{equation*}
Тепер підставимо закон електромагнітної індукції, $\rot\Efield = -\dfrac1c\parttime{\Bfield}$, і матеріальне рівняння $\Bfield=\mu\Hfield$ (так
що $\Hfield\cdot\parttime{\Bfield} = \parttime{}\left(\tfrac12\mu H^2\right)$):
\begin{equation*}
	\vect j\cdot\Efield = -\divg\left(\frac{c}{4\pi}\left[\Efield\times\Hfield\right]\right) - \parttime{}\left(\frac{\epsilon E^2}{8\pi} +
	\frac{\mu H^2}{8\pi}\right).
\end{equation*}

\subsection*{Теорема Пойнтінга}

Проаналізуємо отриманий вираз. Позначимо
\begin{equation}\label{eq:PoyntingVectorDef}
	\tcbhighmath{
		\vect\Pi = \frac{c}{4\pi}\left[\Efield\times\Hfield\right]
	}
\end{equation}
--- \emphx{вектор Пойнтінга}, і $w = \dfrac{\epsilon E^2}{8\pi}+\dfrac{\mu H^2}{8\pi} = \dfrac{\Efield\cdot\Dfield}{8\pi}+\dfrac{\Bfield\cdot\Hfield}{8\pi}$
--- об'ємну густину енергії електромагнітного поля (сума вже відомих нам електричної й магнітної частин, гл.~\ref{ElEnergy},
гл.~\ref{ElMagInduction}). Тоді
\begin{equation}\label{eq:PoyntingTheoremDiff}
	\tcbhighmath{
		-\parttime{w} = \divg\vect\Pi + \vect j\cdot\Efield.
	}
\end{equation}
Інтегруючи по довільному об'єму $V$, обмеженому поверхнею $S$, і застосовуючи теорему Остроградського--Гаусса, отримуємо \emphx{інтегральну}
форму --- \emphx{теорему Пойнтінга}:
\begin{equation}\label{eq:PoyntingTheoremInt}
	\tcbhighmath{
		-\frac{dW}{dt} = \oiint\limits_S \vect\Pi\cdot d\vect S + \iiint\limits_V \vect j\cdot\Efield\, dV.
	}
\end{equation}

\begin{Attention}
	\emphx{Теорема Пойнтінга:} зменшення енергії електромагнітного поля за одиницю часу в даному об'ємі дорівнює сумі потоку енергії крізь
	поверхню, що обмежує цей об'єм, і потужності, яку сили поля виконують над зарядами речовини всередині об'єму. \emphx{Вектор Пойнтінга}
	$\vect\Pi$ визначає кількість енергії поля, що протікає крізь одиничну площадку за одиницю часу, і вказує \emphx{напрям} руху енергії
	електромагнітного поля в просторі --- поняття, яке в статичних задачах електростатики й магнітостатики окремо просто не мало сенсу.
\end{Attention}

\subsection*{Про одиниці вимірювання}

У гауссовій системі одиниць
\begin{equation*}
	\left[\Pi\right] = \frac{\text{ерг}}{\text{см}^2\cdot\text{с}}, \qquad \left[w\right] = \frac{\text{ерг}}{\text{см}^3};
\end{equation*}
для порівняння, у СІ вектор Пойнтінга (там його частіше позначають $\vect S$) вимірюється у Вт/м$^2$, а густина енергії --- у Дж/м$^3$.

\begin{Example}{Pарядка конденсатора}

Плоский конденсатор з пластинами радіуса $R$ заряджається струмом $I$ від зовнішнього джерела (рис.~\ref{pic:CapacitorPoynting}). Якщо в деякий
момент заряд на нижній пластині дорівнює $q$, то вектор індукції всередині конденсатора
\begin{equation}\label{eq:CapacitorD}
	D = 4\pi\sigma = \frac{4\pi q}{\pi R^2},
\end{equation}
напрямлений від нижньої пластини до верхньої. Густина струму зміщення в конденсаторі:
\begin{equation*}
	j_\text{зм} = \frac{\dot D}{4\pi} = \frac{\dot q}{\pi R^2}.
\end{equation*}
Цей струм породжує в конденсаторі магнітне поле, яке знаходимо за теоремою про циркуляцію, взявши за контур $L$ коло радіуса $r$ з центром на
осі конденсатора:
\begin{equation*}
	2\pi r H(r) = \frac{4\pi}{c}\pi r^2 j_\text{зм}(r) \quad\Rightarrow\quad H(r) = \frac{2r}{cR^2}\dot q.
\end{equation*}
Оскільки конденсатор заряджається ($\dot q>0$), струм зміщення напрямлений від нижньої обкладки до верхньої, а силові лінії $\Hfield$
охоплюють вісь конденсатора за правилом правого гвинта (рис.~\ref{pic:CapacitorPoynting}).

\begin{center}%[h!]\centering
	\localinput{capacitor-poynting.tikz}
	\captionof{figure}{Поля $\Efield$, $\Hfield$ і вектор Пойнтінга $\vect\Pi$ при зарядці плоского
		конденсатора\label{pic:CapacitorPoynting}}
\end{center}

Вектор Пойнтінга $\vect\Pi=\dfrac{c}{4\pi}\left[\Efield\times\Hfield\right]$ спрямований \emphx{до осі} конденсатора (перпендикулярно до бічної
поверхні, всередину) і дорівнює за величиною на межі ($r=R$):
\begin{equation*}
	\Pi = \frac{c}{4\pi}EH = \frac{c}{4\pi}\left(\frac{4\pi q}{\epsilon\pi R^2}\right)\left(\frac{2\dot q}{cR}\right) = \frac1{\epsilon\pi
	R^2}\, 2q\dot q.
\end{equation*}
Якщо відстань між пластинами дорівнює $h$, енергія, що втікає в конденсатор за одиницю часу (крізь бічну поверхню площею $2\pi Rh$):
\begin{equation*}
	\frac{dW}{dt} = \Pi\cdot 2\pi Rh = \frac1{\epsilon\pi R^2}\, q\dot q\cdot 2\pi Rh = \frac{q\dot q}{C},
\end{equation*}
де $C=\epsilon\pi R^2/(4\pi h)$ --- ємність конденсатора (гл.~\ref{ElEnergy}). За кінцевий час, коли заряд на пластині досягає значення $q$,
енергія, що втекла в конденсатор ззовні, становить
\begin{equation}\label{eq:CapacitorEnergyResult}
		W = \frac{q^2}{2C}
\end{equation}
--- ми прийшли, зовсім іншим шляхом, до вже відомої формули для енергії зарядженого конденсатора~\eqref{eq:CapacitorEnergy}!
\end{Example}

\begin{Attention}
	Цей результат надзвичайно повчальний: вся енергія, що зрештою запасається в електричному полі \emphx{всередині} конденсатора, насправді
	\emphx{втекла туди ззовні} --- крізь бічну поверхню, а не \emphx{вздовж} проводів чи вздовж силових ліній $\Efield$! Теорема Пойнтінга дає
	абсолютно \emphx{однакову} відповідь на запитання про енергію конденсатора, що й \enquote{наївний} підхід через роботу джерела
	(гл.~\ref{ElEnergy}), --- але картина \emphx{шляху}, яким ця енергія потрапляє в конденсатор, виявляється зовсім іншою й набагато глибшою.
\end{Attention}


\begin{Example}{Довгий провідник}
Розглянемо довгий прямий провід радіуса $R$ і довжини $\ell$, яким тече струм $I$ (рис.~\ref{pic:WirePoynting}). На зовнішній поверхні проводу
присутнє магнітне поле $H = 2I/(cR)$ (гл.~\ref{MagFieldVacuum}). Струм у проводі виникає завдяки напрузі $U$ на його кінцях, що створює всередині
(і, за неперервністю тангенціальної складової, на поверхні) поле $E = U/\ell$.

\begin{center}
	\localinput{wire-poynting.tikz}
	\captionof{figure}{Поля $\Efield$, $\Hfield$ і вектор Пойнтінга $\vect\Pi$ на поверхні довгого
		провідника зі струмом\label{pic:WirePoynting}}
\end{center}

Вектор Пойнтінга виявляється спрямованим \emphx{до осі} провідника й дорівнює за величиною
\begin{equation*}
	\Pi = \frac{c}{4\pi}EH = \frac{c}{4\pi}\cdot\frac{U}{\ell}\cdot\frac{2I}{cR} = \frac1{2\pi\ell R}\,UI.
\end{equation*}
Оскільки площа бічної поверхні провідника дорівнює $S=2\pi\ell R$, повний потік енергії, що втікає в провідник, становить $\Pi S = UI$ --- це в
точності та сама потужність, що йде на джоулеві втрати в проводі: $\Pi S = \mathcal P = UI$ (гл.~\ref{SteadyCurrent}).
\end{Example}

\begin{Attention}
	Цей результат --- один із найбільш контрінтуїтивних, але при цьому й найглибших висновків класичної електродинаміки: \emphx{електрони
		отримують енергію, потрібну для подолання опору й виділення тепла, не \enquote{вздовж} проводу від джерела, а ззовні, з поля, що оточує
		провід}! Енергія, віддана далеким джерелом (батареєю), якимось чином \enquote{розтікається} через навколишній простір і \emphx{втікає}
	всередину провідника крізь його бічну поверхню, а сам провід відіграє роль лише напрямної труби для цього потоку енергії, а не
	\enquote{кабелю}, яким енергія тече вздовж. Це разюче суперечить наївному (але дуже поширеному) уявленню, ніби енергія \enquote{тече по
		дротах} так само, як вода тече по трубі.
\end{Attention}

%% --------------------------------------------------------
\section{Скін-ефект}
%% --------------------------------------------------------

До цього моменту рівняння Максвелла поставали перед нами або як \emphx{підсумок} уже відомого (стаціонарні випадки --- електро- і
магнітостатика), або як \emphx{джерело} якісних висновків (струм зміщення, симетрія, теорема Пойнтінга). Тепер подивимося на них як на
\emphx{робочий інструмент}, узявши конкретну фізичну задачу: поведінку електромагнітного поля у провіднику за швидкозмінних процесів, на прикладі скін-ефекту.

%% --------------------------------------------------------
\subsection*{Означення та наслідки}
%% --------------------------------------------------------

\emphx{Скін-ефектом} (від англ. \emph{skin} --- шкіра) називають протікання струмів високої частоти переважно в тонкому поверхневому шарі
провідника (рис.~\ref{pic:SkinIllustration}).

\begin{figure}[h!]\centering
	\localinput{skin.tikz}
	\caption{Зменшення ефективної товщини струмонесучого шару зі зростанням
		частоти\label{pic:SkinIllustration}}
\end{figure}

Скін-ефект має кілька важливих практичних наслідків:
\begin{itemize}
	\item зменшується \emphx{діюча} площа перерізу провідника, а отже, \emphx{зростає} його опір на змінному струмі;
	\item \emphx{зменшується} індуктивність провідника: оскільки за сильного скін-ефекту струму (а отже, і магнітного поля) всередині
	      провідника практично немає, магнітна енергія стає меншою на величину енергії поля, яка раніше запасалась усередині провідника;
	\item електромагнітна хвиля, що потрапляє на поверхню провідника (металу, електроліту, плазми), швидко загасає вглиб нього, проникаючи лише
	      на невелику глибину --- саме тому \emphx{провідники добре відбивають світло}.
\end{itemize}

\subsection*{Якісне пояснення}

Змінний струм у провіднику породжує змінне вихрове магнітне поле, силові лінії якого перпендикулярні до осі провідника
(рис.~\ref{pic:SkinEffectQualitative}). Це змінне магнітне поле, своєю чергою, породжує вихрове електричне поле (гл.~\ref{ElMagInduction}), що
спричиняє в товщі провідника струми Фуко: на \emphx{поверхні} провідника вихровий струм напрямлений \emphx{у той самий бік}, що й основний струм
провідності (і \emphx{підсилює} його там), а \emphx{всередині} провідника --- у \emphx{протилежний} бік (і \emphx{послаблює} струм у серцевині).
У результаті струм \emphx{витісняється} з осі провідника до його поверхні.

\begin{SCfigure}[3][h!]\centering
	\localinput{skin-effect.tikz}
	\caption{Якісний механізм скін-ефекту: вихрові струми Фуко підсилюють струм на поверхні й
		послаблюють у центрі провідника\label{pic:SkinEffectQualitative}}
\end{SCfigure}

\subsection*{Чому в металах нехтують струмами зміщення}

У провіднику густина струму провідності й напруженість поля пов'язані законом Ома, $\vect j = \lambda\Efield$. Якщо поле в провіднику змінюється
гармонічно, $\Efield=\Efield_0\cos(\omega t)$, то поряд зі струмом провідності виникає й струм зміщення:
\begin{equation*}
	\vect j_\text{зм} = \frac{\epsilon}{4\pi}\parttime{\Efield} = \frac{\epsilon\omega}{4\pi}E_0\cos\left(\omega t+\frac\pi2\right),
\end{equation*}
величину $\dfrac{\epsilon\omega}{4\pi}$ можна формально трактувати як \enquote{питому провідність} для струмів зміщення. Для типових металів
провідність $\lambda = 10^{15}$--$10^{17}$~с$^{-1}$ (гауссова система), а типова діелектрична проникність металів --- порядку $10$--$100$; тому аж
до частот порядку $10^{13}$~Гц виконується нерівність
\begin{equation}\label{eq:SkinDisplacementNegligible}
	\frac{\epsilon\omega}{4\pi} \ll \lambda,
\end{equation}
тобто у широкому діапазоні частот струми зміщення \emphx{не відіграють істотної ролі} у металах порівняно зі струмами провідності --- саме тому
надалі, аналізуючи скін-ефект у металах, ми зможемо ними знехтувати.

%% --------------------------------------------------------
\subsection*{Строге виведення глибини проникнення}
%% --------------------------------------------------------

Виведемо глибину проникнення поля в провідник, спираючись безпосередньо на рівняння Максвелла (а не на якісні міркування). Нехай $\epsilon=\const$,
$\mu=\const$; у провіднику немає некомпенсованого об'ємного заряду, тому теорема Гаусса дає $\divg\Efield=0$. З урахуванням закону Ома, а також
знехтувавши, за щойно встановленою нерівністю~\eqref{eq:SkinDisplacementNegligible}, струмом зміщення, рівняння Максвелла в металі набувають
вигляду
\begin{equation*}
	\rot\Efield = -\frac{\mu}c\parttime{\Hfield}, \qquad \rot\Hfield = \frac{4\pi\lambda}{c}\Efield, \qquad \divg\Efield=0, \quad \divg\Hfield=0.
\end{equation*}
Виключимо $\Hfield$: застосуємо операцію $\rot$ до першого рівняння й підставимо друге:
\begin{equation*}
	\rot\rot\Efield = -\frac{\mu}c\parttime{}\left(\rot\Hfield\right) = -\frac{4\pi\mu\lambda}{c^2}\parttime{\Efield}.
\end{equation*}
З іншого боку, за тотожністю векторного аналізу і з урахуванням $\divg\Efield=0$: $\rot\rot\Efield = \grad(\divg\Efield)-\nabla^2\Efield =
-\nabla^2\Efield$. Прирівнюючи, отримуємо \emphx{рівняння скін-ефекту}:
\begin{equation}\label{eq:SkinWaveEquation}
	\tcbhighmath{
		\nabla^2\Efield = \frac{4\pi\mu\lambda}{c^2}\parttime{\Efield}.
	}
\end{equation}

%% --------------------------------------------------------
\subsection*{Оцінка глибини проникнення}
%% --------------------------------------------------------

Рівняння~\eqref{eq:SkinWaveEquation} за математичною формою збігається з рівнянням \emphx{дифузії}, $\nabla^2 F = D\parttime{F}$, з коефіцієнтом
дифузії $D = c^2/(4\pi\mu\lambda)$. Якщо на поверхні речовини діє джерело, то за час $\tau$ \enquote{хвиля дифузії} проникає на глибину
$\delta=\sqrt{D\tau}$; якщо джерело діє періодично з періодом $T=2\pi/\omega$, слід узяти $\tau\approx T$, звідки отримуємо \emphx{оцінку} для
глибини проникнення:
\begin{equation}\label{eq:SkinDepthEstimate}
	\delta \sim \frac{c}{\sqrt{4\pi\lambda\mu\omega}}.
\end{equation}

Переконаємось, що ця оцінка, з точністю до числового множника, дійсно правильна, застосувавши безпосередньо рівняння Максвелла в інтегральній
формі до тонкого приповерхневого шару завтовшки $\delta$.

\emphx{Циркуляція $\Efield$.} Розглянемо прямокутний контур $L$ висотою $h$ і завширшки $\delta$ уздовж осі провідника
(рис.~\ref{pic:SkinEstimate1}). За законом електромагнітної індукції, $Eh = -\dfrac1c\dfrac{\Delta\Phi}{\Delta t}$, де магнітний потік крізь цей
контур $\Phi = -BS = -\mu H h\delta$ (поле $H$ спадає до нуля на глибині $\delta$), а швидкість його зміни за порядком величини $\Delta\Phi/\Delta
t \approx \omega\Phi$ ($\omega\sim1/T$ --- частота зміни струму). Тому
\begin{equation}\label{eq:SkinCirculationE}
	Eh \sim \frac1c\mu H\omega h\delta \quad\Rightarrow\quad E \sim \frac{\mu H\omega\delta}{c}.
\end{equation}

%\begin{figure}[h!]\centering
%	\localinput{skin-estimate-1.tikz}
%	\caption{Контур для оцінки циркуляції $\Efield$ поблизу поверхні провідника\label{pic:SkinEstimate1}}
%\end{figure}

\emphx{Циркуляція $\Hfield$.} Тепер розглянемо контур $L$ на поперечному перерізі провідника, завдовжки $\ell$ уздовж поверхні й завширшки
$\delta$ углиб (рис.~\ref{pic:SkinEstimate2}). За теоремою про циркуляцію, $H\ell = \dfrac{4\pi}{c}j\ell\delta$, а з огляду на закон Ома,
$j=\lambda E$:
\begin{equation}\label{eq:SkinCirculationH}
	H \sim \frac{4\pi\lambda}{c}E\delta.
\end{equation}

\begin{figure}[h!]
	\centering

	\captionbox{Контур для оцінки циркуляції $\Efield$ поблизу поверхні провідника\label{pic:SkinEstimate1}}[0.47\textwidth]{%
		\localinput{skin-estimate-1.tikz}%
	}
	\hfill % Розсуває малюнки по краях сторінки
	\captionbox{Контур для оцінки циркуляції $\Hfield$ поблизу поверхні провідника\label{pic:SkinEstimate2}}[0.47\textwidth]{%
		\localinput{skin-estimate-2.tikz}%
	}
\end{figure}


%\begin{figure}[h!]\centering
%	\localinput{skin-estimate-2.tikz}
%	\caption{Контур для оцінки циркуляції $\Hfield$ на перерізі провідника\label{pic:SkinEstimate2}}
%\end{figure}

Підставляючи~\eqref{eq:SkinCirculationH} в~\eqref{eq:SkinCirculationE}:
\begin{equation*}
	E \sim \frac{4\pi\mu\lambda\omega\delta^2}{c^2}\, E,
\end{equation*}
звідки, скорочуючи на $E$, отримуємо \emphx{глибину скін-шару}:
\begin{equation}\label{eq:SkinDepthFinal}
	\tcbhighmath{
		\delta = \frac{c}{\sqrt{4\pi\lambda\mu\omega}}
	}
\end{equation}
--- той самий результат, що й дала якісна оцінка через аналогію з дифузією~\eqref{eq:SkinDepthEstimate}, тепер уже виведений із перших принципів
(з точністю до числового множника порядку одиниці, який строгий розв'язок хвильового рівняння~\eqref{eq:SkinWaveEquation} підтверджує рівним
точно $1/\sqrt{4\pi}$).

\begin{Attention}
	Отримана формула для $\delta$ --- це результат оцінки \enquote{за порядком величини}, побудованої на спрощених контурних міркуваннях, а не
	точний математичний розв'язок. Числовий коефіцієнт $1/\sqrt{4\pi}$ у ній --- радше щасливий збіг оцінки з точним результатом, ніж
	закономірність: строгий аналіз, що спирається безпосередньо на розв'язання хвильового рівняння~\eqref{eq:SkinWaveEquation} (комплексним
	методом, з експоненційно загасаючим розв'язком $E\propto e^{-x/\delta}\cos(\omega t - x/\delta)$), підтверджує цей самий результат із високою
	точністю й для різних геометрій. Але в загальному випадку оцінковий підхід не \emphx{гарантує} правильного числового коефіцієнта --- лише
	правильний порядок величини й правильну залежність від параметрів задачі ($\delta\propto 1/\sqrt{\lambda\mu\omega}$), і збіг тут варто
	сприймати саме так.
\end{Attention}

\begin{table}[h!]\small\centering
	\caption{Глибина скін-шару для деяких металів на частоті $10$~ГГц}
		\begin{tblr}{
		colspec={X[l, m]X[c,m]},
        row{odd} = {gray!10},
		row{1} = {c, m, fg=white, bg=themecolorlight, font=\bfseries},
		hlines,
		vlines
		}
		Метал         & Глибина скін-шару, мкм \\
		Мідь (\ce{Cu})   & 0.652                    \\
		Срібло (\ce{Ag}) & 0.634                    \\
		Золото (\ce{Au}) & 0.753                    \\
		Алюміній (\ce{Al})& 0.652                   \\
	\end{tblr}
\end{table}

%% --------------------------------------------------------
\subsection*{Застосування скін-ефекту в техніці}
%% --------------------------------------------------------

Існування скін-ефекту завжди враховують у техніці швидкозмінних струмів. Оскільки такі струми практично не йдуть у глибині провідника, лінії
електропередач для них \emphx{виготовляють із порожнистих труб} --- економлячи метал без істотної втрати в опорі. У радіотехніці надвисоких частот
багато деталей (хвилеводи, коаксіальні лінії) покривають тонким, добре провідним шаром срібла, оскільки їхній опір практично повністю
визначається лише поверхневим шаром. Скін-ефект у поєднанні з нагріванням вихровими струмами (гл.~\ref{ElMagInduction}) використовують і для
\emphx{поверхневого загартування сталі}: деталі піддають дії змінних струмів і потім швидко охолоджують, унаслідок чого розігрів і загартування
відбуваються лише в поверхневому шарі, а основна маса деталі лишається незагартованою --- це дає вироби з високою стійкістю до стирання поверхні,
але без крихкості, властивої звичайному наскрізному загартуванню.

%% --------------------------------------------------------
\section{Про рівняння Максвелла}
%% --------------------------------------------------------

\emphx{Теорія Максвелла --- це вершина класичної науки про електромагнітне поле}. До цієї висоти ми піднімалися крок за кроком, уточнюючи поняття
й розкриваючи глибинні принципи цієї науки: кожен пройдений розділ --- електростатика, закони постійного струму, магнітостатика --- був немов
базовий табір на шляху до вершини.

\begin{SCfigure}[4][h!]\centering
	\includegraphics[height=4cm]{Hertz}
	\caption{Генріх Рудольф Герц (1857--1894) --- перший експериментально отримав
		електромагнітні хвилі, теоретично передбачені Максвеллом\label{pic:Hertz}}
\end{SCfigure}

На вершині перед нами постають рівняння Максвелла --- велична симфонія електричного й магнітного полів, ключ до розуміння електромагнітної
взаємодії. Але шлях науки триває: з вершини ми спускаємось у простір застосувань --- досліджуємо поширення електромагнітних хвиль, принципи
випромінювання й закладаємо основи оптики, радіо- й електротехніки. Ця подорож --- не лише тріумф інтелекту, а й людського духу, що прагне
підкорювати вершини й відкривати нові горизонти.



\begin{tornpage}\small\itshape
	Неможливо позбутися відчуття, що математичні формули живуть власним життям і володіють власним розумом, що вони розумніші за нас, розумніші
	навіть за тих, хто їх відкрив, що ми отримуємо з цих формул більше, ніж у них було спочатку закладено\ldots\\

   \hfill{Г. Герц}.
\end{tornpage}

%\begin{figure}[h!]\centering
%	\includegraphics[height=3.5cm]{ontop}
%	\caption{На вершині класичної електродинаміки\label{pic:OnTop}}
%\end{figure}

Уся споруда класичної електродинаміки, зведена нами розділ за розділом --- від закону Кулона до струму зміщення, --- вміщується тепер у
чотирьох компактних рівняннях Максвелла, доповнених трьома матеріальними рівняннями середовища. Ці чотири рівняння містять \emphx{усе}:
електростатику й магнітостатику як окремі стаціонарні випадки, закон Ома як наслідок для провідних середовищ, електромагнітну індукцію й
струм зміщення як дві половини єдиної симетрії, --- і, як буде показано далі, \emphx{електромагнітні хвилі} й саме \emphx{світло} як прямий
наслідок цих же рівнянь у вакуумі. Саме в цьому й полягає та \enquote{глибина, що виявляється тільки при ретельному вивченні}, про яку писали
Ейнштейн та Інфельд: прості на вигляд чотири рівняння виявляються не описом окремого явища, а повною структурою одного з фундаментальних
полів природи.


\begin{Questions}
\begin{quizlist}
	\item У чому полягала ідея Максвелла про симетрію електричного й магнітного полів? Як ця ідея привела його до поняття струму зміщення?
	\item Покажіть, що теорема про циркуляцію $\Hfield$ для постійних полів, застосована до випадку змінного в часі заряду, суперечить закону збереження заряду (рівнянню неперервності).
	\item Виведіть вираз для густини струму зміщення $\vect j_\text{зм} = \dfrac1{4\pi}\parttime{\Dfield}$ з вимоги узгодженості теореми про циркуляцію із законом збереження заряду. Чому цей вибір не єдиний математично можливий, і що зрештою підтверджує саме його правильність?
	\item Запишіть остаточну теорему про циркуляцію магнітного поля з урахуванням струму зміщення в інтегральній та диференціальній формах.
	\item Розгляньте задачу про радіальне стікання заряду з кулі: чому магнітне поле в цій задачі не виникає взагалі, попри те що заряд явно рухається? Яку роль тут відіграє точна компенсація струму провідності струмом зміщення?
	\item Розгляньте задачу про зарядку плоского конденсатора: чому циркуляція $\Hfield$ по одному й тому самому контуру дає однаковий результат для двох різних поверхонь --- тієї, що перетинає провід, і тієї, що проходить між обкладками? Що це означає для ролі струму зміщення в конденсаторі?
	\item Чим струм зміщення принципово відрізняється від струму провідності --- чому до нього незастосовні закон Ома і закон Джоуля--Ленца?
	\item Порівняйте закон електромагнітної індукції Фарадея та гіпотезу Максвелла про магнітоелектричну індукцію. У чому полягає асиметрія знаків між ними і до яких глибоких наслідків вона приводить?
	\item Запишіть повну систему чотирьох рівнянь Максвелла в інтегральній та диференціальній формах. За яким принципом ці рівняння групуються в дві пари?
	\item Виведіть вираз електричного поля через потенціали, $\Efield = -\grad\phi - \dfrac1c\parttime{\vect A}$, з теореми Гаусса для $\Bfield$ і закону індукції Фарадея. Чому в нестаціонарному випадку $\Efield$ визначається не лише скалярним потенціалом?
	\item Покажіть, що в стаціонарному випадку повна система рівнянь Максвелла розпадається на незалежні системи електростатики й магнітостатики. Чому зв'язок електрики й магнетизму --- принципово нестаціонарне явище?
	\item Перелічіть матеріальні рівняння, що доповнюють рівняння Максвелла ($\Dfield=\epsilon\Efield$, $\Bfield=\mu\Hfield$, закон Ома). Навіщо вони потрібні для розв'язання конкретних задач?
	\item Запишіть граничні умови для електричних і магнітних векторів на межі розділу двох середовищ у загальному (нестаціонарному) випадку. Чим вони відрізняються від умов для стаціонарних полів?
	\item Виведіть теорему Пойнтінга $-\parttime{w} = \divg\vect\Pi + \vect j\cdot\Efield$ з рівнянь Максвелла для $\rot\Hfield$ і $\rot\Efield$, використовуючи тотожність для дивергенції векторного добутку.
	\item Що таке вектор Пойнтінга і густина енергії електромагнітного поля? Сформулюйте теорему Пойнтінга в інтегральній формі та поясніть фізичний зміст кожного доданка.
	\item Розгляньте приклад зарядки конденсатора через теорему Пойнтінга: чому вся енергія, що запасається в полі всередині конденсатора, насправді втікає туди ззовні, крізь бічну поверхню, а не вздовж проводів?
	\item Розгляньте приклад довгого провідника зі струмом: чому енергія, що йде на джоулеві втрати, надходить у провідник не вздовж нього від джерела, а ззовні, крізь бічну поверхню?
	\item Що таке скін-ефект і які його основні практичні наслідки (зміна опору, індуктивності, відбивання світла провідниками)?
	\item Дайте якісне пояснення скін-ефекту через вихрові струми Фуко: чому вони підсилюють струм на поверхні провідника й послаблюють у його серцевині?
	\item Чому в металах аж до дуже високих частот ($\sim10^{13}$ Гц) можна нехтувати струмом зміщення порівняно зі струмом провідності? Виведіть умову $\epsilon\omega/(4\pi) \ll \lambda$.
	\item Виведіть рівняння скін-ефекту $\nabla^2\Efield = (4\pi\mu\lambda/c^2)\,\parttime{\Efield}$ з рівнянь Максвелла для провідника, знехтувавши струмом зміщення.
	\item Виведіть оцінку глибини скін-шару $\delta \sim c/\sqrt{4\pi\lambda\mu\omega}$, спираючись на аналогію з рівнянням дифузії, а також прямим контурним міркуванням через циркуляції $\Efield$ і $\Hfield$ біля поверхні провідника.
	\item Які практичні застосування скін-ефекту в техніці (порожнисті провідники ЛЕП, посріблення хвилеводів, поверхневе загартування сталі)?
\end{quizlist}
\end{Questions}
